\chapter{Robotics}

\section{Embodiment Changes the Constraint, Not the Architecture}

Chapter 32 asked what this architecture buys a simulation a scientist observes from outside it. This chapter asks about a harder case: a system whose proposals must be realized as actual physical motion, by actual hardware, subject to actual torque limits, joint ranges, and the unforgiving fact that a physically unrealizable \(\Delta\psi\) cannot simply be logged as an interesting hypothesis and set aside. Embodiment does not require a new architecture. It requires a new member of C. Physical realizability, whether a proposed change to the robot's own joint configuration corresponds to something the actual actuators can execute within their real limits, joins the constraint manifold exactly as Chapter 32's conservation laws did, a hard boundary \(\Pi_C\) enforces rather than a preference the controller has merely learned to respect.

\section{The Robot as Observer, Generator, and Body}

A robot occupies all three roles Chapter 6 distinguished, observer, generator, and world, in an unusually tight loop. It observes, through sensors realizing \(O_\theta\). It generates, proposing actions as \(\Delta\psi\) through exactly the interaction machinery Chapter 28 built. And its own body, joint angles, limb positions, current gait phase, is itself part of \(M_t\), not external to the field it is acting on. This is the technical content behind the word embodiment in this framework: a robot is not an agent standing outside a field and acting on it from a safe remove, the way a human player's avatar was implicitly treated in earlier chapters. It is a region of the field proposing changes to itself, and to the field immediately around it, simultaneously.

\section{Central Pattern Generators as Dynamic Fixed Points}

Locomotion supplies this chapter's central example, and it is worth stating its core claim plainly before developing it: gait is not a sequence of target poses a controller repeatedly corrects toward. It is a dynamic fixed point in exactly Chapter 16's sense, later refined by Chapter 20 into a conserved trajectory through phase space. A central pattern generator, the biological and robotic mechanism that produces rhythmic locomotor patterns through coupled oscillation rather than explicit trajectory planning, was already trajectory-centric before this book's own vocabulary existed to describe it this way; this chapter's contribution is recognizing that a walking gait is precisely the kind of continuously regenerated coherence Chapter 16 built its account of dynamic fixed points to describe, not the kind of static, trivially-stable target pose a naive controller might instead pursue. Balance, under this reading, is not a state a robot achieves and then holds still. It is a coherence, Φ, continually regenerated by an ongoing cycle of proposal, projection, and commit, in the same sense the fountain from Chapter 16 was never at rest even while its overall form remained recognizably stable.

\section{Hierarchical Composition and the Kuramoto Coupling}

A single joint's oscillation is described by a phase, and a population of coupled joint-level oscillators synchronizes according to a coupling law most naturally expressed through the Kuramoto order parameter,

\[ Z(t) = r(t)\, e^{i\psi(t)}, \]

where r(t) measures how synchronized the population currently is and \(\psi(t)\) is the population's mean phase. This is Chapter 20's phase-space machinery, Φ and v as conjugate pairs generating conservative oscillation, applied to an unusually literal case: a joint's own oscillatory state genuinely is a position-and-momentum-like pair in the sense Chapter 20 developed, not merely analogous to one. What makes locomotion architecturally interesting, and what this chapter borrows directly from existing work on hierarchical continuation in gait, is that this coupling does not stop at a single joint. Joint-level oscillator populations compose into limb-level coordination; limb-level operators compose into gait-level pattern, walking, trotting, galloping; and gait-level pattern composes into whole-body coordination. Different gaits, under this composition, are not different dynamical laws requiring separate controllers. They are different parameterizations of the same underlying coupling law, distinguished only by different target phase offsets in the coupling matrix K relating one oscillator population to another. This is Chapter 23's structural-equivalence question, answered concretely: a walk and a trot are isomorphic in the relevant sense, related by a structure-preserving reparameterization of one shared coupling architecture, not two unrelated things that happen to both move a robot forward.

Each level of this hierarchy also runs on its own clock, in precisely Chapter 17's sense. Individual joint oscillation is the fastest layer, cycling many times within a single stride. Limb coordination is slower. Overall gait pattern slower still. Whole-body balance, the slowest layer, is what Chapter 25.2's scheduler would treat as this hierarchy's coarsest update rate, checked and reconfirmed far less often than any individual joint's phase, exactly as Chapter 17 argued a field's slow components generally are.

\section{A Sketch of the CPG Chain Controller}

It is worth making this hierarchy's shape concrete before moving on, both as a controller design and as an instance of Chapter 25's system architecture. The composition runs as follows:

\begin{Verbatim}[breaklines=true, breaksymbolleft={}, fontsize=\small]
JointOscillator
  - phase: real
  - naturalFrequency: real
  - couples_to: [JointOscillator]  (via coupling weight K_ij)
  + step(dt) -> proposes ΔΦ_joint, Δv_joint

LimbController
  - oscillators: [JointOscillator]
  - orderParameter: Z(t) = r(t) e^{i ψ(t)}
  + synchronize() -> aggregates joint proposals into one limb-level Δψ

GaitController
  - limbs: [LimbController]
  - couplingMatrix: K            (target phase offsets; walk/trot/gallop
                                   differ only in K's values, not its shape)
  - activeGait: GaitParameterization
  + selectGait(K) -> reparameterizes limb coupling, no controller swap

WholeBodyCoordinator
  - gait: GaitController
  - balanceCoherence: Φ_whole_body
  + monitor() -> flags Φ_whole_body decline before observable instability

ProposalService (Chapter 25)
  <- receives Δψ from WholeBodyCoordinator, same interface as any
     other proposal source; no special-cased path for "robot" proposals
\end{Verbatim}

Nothing in this sketch introduces a component Chapter 25 did not already specify. WholeBodyCoordinator's output enters the proposal service exactly as any other \(\Delta\psi\) would, and is projected against C, now including the physical-realizability constraint from section 33.1, exactly as any other proposal is. A CPG chain controller, under this architecture, is not a special robotics subsystem sitting beside the world engine. It is a particular, highly structured way of generating \(\Delta\psi\), nothing more privileged than that.

\section{A Grammar for CPG Chain Configuration}

Because different gaits are reparameterizations of one coupling law rather than separate controllers, it is useful to have a small, well-formed grammar for specifying a configuration rather than hand-coding each gait separately. The following BNF sketch is deliberately minimal, sufficient to specify a coupling matrix and its target phase offsets rather than a complete robotics description language:

\begin{Verbatim}[breaklines=true, breaksymbolleft={}, fontsize=\small]
<cpg-config>      ::= <oscillator-list> <coupling-list> <gait-name>

<oscillator-list>  ::= "oscillators" "{" <oscillator> { "," <oscillator> } "}"
<oscillator>       ::= <id> ":" "freq" "=" <real>

<coupling-list>    ::= "coupling" "{" <coupling-entry> { "," <coupling-entry> } "}"
<coupling-entry>   ::= <id> "->" <id> ":" "weight" "=" <real>
                                    "," "phase-offset" "=" <real>

<gait-name>        ::= "gait" "=" <string>

<id>               ::= letter { letter | digit }
<real>              ::= digit { digit } [ "." digit { digit } ]
<string>            ::= "\"" { character } "\""
\end{Verbatim}

A walking gait and a trotting gait, under this grammar, are two instances differing only in their `<coupling-list>`'s phase-offset values and the `<gait-name>` label; the `<oscillator-list>` and the shape of the coupling relation can remain identical. This gives concrete, checkable content to the claim from section 33.4 that gaits are parameterizations rather than separate designs: two configurations sharing an oscillator list and coupling topology, differing only in phase offsets, are exactly the isomorphic-but-differently-parameterized objects Chapter 23 gives a name to.

\section{Balance as a Diagnostic: Coherence Before Collapse}

This architecture offers one further, genuinely practical payoff worth naming. Because whole-body coherence, Φ, is tracked explicitly rather than inferred only from whether the robot has actually fallen, it becomes possible to monitor coherence as a leading indicator rather than waiting for failure to become visible. A stumble, on this reading, is a trajectory moving from the neighborhood of a stable fixed point, Chapter 19.5's local minimum of the relevant energy functional, toward an unstable one, a saddle or local maximum from which small perturbations grow rather than being absorbed back. This book's own machinery predicts that Φ should measurably decline during this transition before any externally observable loss of balance occurs, since the fixed point's stability, not merely the robot's current pose, is what is actually degrading. A controller monitoring Φ directly, rather than only monitoring whether the robot is still upright, has a genuine early-warning signal this architecture supplies for free, simply by taking its own account of coherence and stability seriously as something worth measuring in real time rather than only as a theoretical construct.

\section{What This Chapter Has Not Solved}

This chapter has not designed a coupling matrix, tuned a controller for any particular robot body, or specified how a physical-realizability constraint's actual torque and range limits should be derived from a given robot's hardware. Those are control-theoretic and mechanical engineering questions, exactly as the equations governing a scientific simulation were, in Chapter 32, the domain scientist's responsibility rather than this book's own contribution. What this chapter has offered is a place for an existing, independently developed body of locomotor control theory to sit inside this book's architecture without contradiction, and a small amount of genuine payoff, the coherence-decline diagnostic in particular, that falls out of taking the architecture's own vocabulary seriously.

\section{Looking Ahead}

Robotics asked what happens when a persistent field must be embodied in hardware that acts on and within the world it represents. Chapter 34 turns to a different kind of embodiment, not in physical machinery but in a learner's own developing understanding, asking what this architecture offers education: a domain where the field being maintained is not a robot's body or a physical system, but a student's evolving grasp of one.
